% GENERATED FILE. Edit canonical lessons, then run npm run generate:books.
% canonical-source-hash: fnv1a64:d8b1cb4a85500d2f
% canonical-lessons: AR-C470-yanbu, AR-C470-manba, AR-C470-masabb, AR-C470-diffa, AR-C470-adhb, AR-R470-review-463-470

\chapter{Springs and rivers}
\label{ch:ar-springs-and-rivers}
\hlchaptermodality{470}

\begin{chapteropening}
\textbf{By the end of this chapter:} \emph{I can say a spring, a river source, a river mouth, a riverbank and sweet water.}

Complete the last lesson of chapter 470: I can say a spring, a river source, a river mouth, a riverbank and sweet water.
\end{chapteropening}

\section[\texorpdfstring{\ar{ينبوع} (yanbūʿ)}{ينبوع (yanbūʿ)}]{\ar{ينبوع} (yanbūʿ) — a spring (of water)}

\label{lesson:AR-C470-yanbu}



\begin{quote}
Before the new one: say the Arabic for fertile, then the Arabic for to uproot.
\end{quote}

\subsection*{You'll want to know: \ar{ينبوع}}
\textbf{\ar{ينبوع}} — \emph{yanbūʿ} — \textquotedblleft{}a spring (of water)\textquotedblright{}.

\begin{quote}
\ar{شربنا} \ar{من} \ar{ينبوع} \ar{في} \ar{الجبل}.
\end{quote}

\begin{quote}
\emph{sharibnā min yanbūʿ fī al-jabal.}
\end{quote}

\begin{quote}
We drank from a spring in the mountain.
\end{quote}

Plural: \ar{ينابيع}, \emph{yanābīʿ}.

\begin{cousinweb}
\textbf{\ar{ينبوع}} (\emph{yanbūʿ}) is built on the native Arabic root \emph{n-b-ʿ}, \textquotedblleft{}to well up, to gush\textquotedblright{} (Lane). EtymDB traces the word back to Proto-Semitic.
\end{cousinweb}

\begin{grammarlens}[title={What you've built}]
One more word for getting around.
\end{grammarlens}

\subsection*{Your turn}
\begin{itemize}
\raggedright
  \item \emph{Say it:} \emph{yanbūʿ}
  \item \emph{Say it:} \emph{yanbūʿ}, once more
  \item \emph{Say it:} \emph{sharibnā min yanbūʿ fī al-jabal.} — We drank from a spring in the mountain.
  \item \emph{Recall:} say the Arabic for fertile, then the Arabic for to uproot, then say \emph{yanbūʿ} again
\end{itemize}

\begin{tcolorbox}[breakable,colback=teal!4,colframe=teal!35!black,title={Before you move on}]
What does \ar{ينبوع} mean? (\textquotedblleft{}A spring (of water)\textquotedblright{}.) Say it once more.
\end{tcolorbox}
\section[\texorpdfstring{\ar{منبع} (manbaʿ)}{منبع (manbaʿ)}]{\ar{منبع} (manbaʿ) — a river source; a source}

\label{lesson:AR-C470-manba}



\begin{quote}
Before the new one: say the Arabic for to uproot, then the Arabic for a spring.
\end{quote}

\subsection*{You'll want to know: \ar{منبع}}
\textbf{\ar{منبع}} — \emph{manbaʿ} — \textquotedblleft{}a river source; a source\textquotedblright{}.

\begin{quote}
\ar{يقع} \ar{منبع} \ar{النهر} \ar{في} \ar{الجبال}.
\end{quote}

\begin{quote}
\emph{yaqaʿu manbaʿ an-nahr fī al-jibāl.}
\end{quote}

\begin{quote}
The source of the river lies in the mountains.
\end{quote}

\begin{cousinweb}
\textbf{\ar{منبع}} (\emph{manbaʿ}) is built on the native Arabic root \emph{n-b-ʿ}, \textquotedblleft{}to well up, to gush\textquotedblright{} (Lane). The pattern \emph{mafʿal} makes it the place where water wells up.
\end{cousinweb}

\begin{grammarlens}[title={What you've built}]
One more word for getting around.
\end{grammarlens}

\subsection*{Your turn}
\begin{itemize}
\raggedright
  \item \emph{Say it:} \emph{manbaʿ}
  \item \emph{Say it:} \emph{manbaʿ}, once more
  \item \emph{Say it:} \emph{yaqaʿu manbaʿ an-nahr fī al-jibāl.} — The source of the river lies in the mountains.
  \item \emph{Recall:} say the Arabic for to uproot, then the Arabic for a spring, then say \emph{manbaʿ} again
\end{itemize}

\begin{tcolorbox}[breakable,colback=teal!4,colframe=teal!35!black,title={Before you move on}]
What does \ar{منبع} mean? (\textquotedblleft{}A river source; a source\textquotedblright{}.) Say it once more.
\end{tcolorbox}
\section[\texorpdfstring{\ar{مصب} (maṣabb)}{مصب (maṣabb)}]{\ar{مصب} (maṣabb) — a river mouth}

\label{lesson:AR-C470-masabb}



\begin{quote}
Before the new one: say the Arabic for a spring, then the Arabic for a river source.
\end{quote}

\subsection*{You'll want to know: \ar{مصب}}
\textbf{\ar{مصب}} — \emph{maṣabb} — \textquotedblleft{}a river mouth\textquotedblright{}.

\begin{quote}
\ar{عند} \ar{المصب} \ar{يلتقي} \ar{النهر} \ar{بالبحر}.
\end{quote}

\begin{quote}
\emph{ʿinda al-maṣabb yaltaqī an-nahr bi-l-baḥr.}
\end{quote}

\begin{quote}
At the mouth the river meets the sea.
\end{quote}

\begin{cousinweb}
\textbf{\ar{مصب}} (\emph{maṣabb}) is built on the native Arabic root \emph{ṣ-b-b}, \textquotedblleft{}to pour\textquotedblright{} (Lane). It is the place where a river pours out; \emph{ṣabba}, \textquotedblleft{}to pour\textquotedblright{}, shares the root.
\end{cousinweb}

\begin{grammarlens}[title={What you've built}]
One more word for getting around.
\end{grammarlens}

\subsection*{Your turn}
\begin{itemize}
\raggedright
  \item \emph{Say it:} \emph{maṣabb}
  \item \emph{Say it:} \emph{maṣabb}, once more
  \item \emph{Say it:} \emph{ʿinda al-maṣabb yaltaqī an-nahr bi-l-baḥr.} — At the mouth the river meets the sea.
  \item \emph{Recall:} say the Arabic for a spring, then the Arabic for a river source, then say \emph{maṣabb} again
\end{itemize}

\begin{tcolorbox}[breakable,colback=teal!4,colframe=teal!35!black,title={Before you move on}]
What does \ar{مصب} mean? (\textquotedblleft{}A river mouth\textquotedblright{}.) Say it once more.
\end{tcolorbox}
\section[\texorpdfstring{\ar{ضفة} (ḍiffa)}{ضفة (ḍiffa)}]{\ar{ضفة} (ḍiffa) — a riverbank}

\label{lesson:AR-C470-diffa}



\begin{quote}
Before the new one: say the Arabic for a river source, then the Arabic for a river mouth.
\end{quote}

\subsection*{You'll want to know: \ar{ضفة}}
\textbf{\ar{ضفة}} — \emph{ḍiffa} — \textquotedblleft{}a riverbank\textquotedblright{}.

\begin{quote}
\ar{جلسنا} \ar{على} \ar{ضفة} \ar{النهر}.
\end{quote}

\begin{quote}
\emph{jalasnā ʿalā ḍiffat an-nahr.}
\end{quote}

\begin{quote}
We sat on the riverbank.
\end{quote}

Plural: \ar{ضفاف}, \emph{ḍifāf}.

\begin{cousinweb}
\textbf{\ar{ضفة}} (\emph{ḍiffa}) is usually traced to the native Arabic root \emph{ḍ-f-f}, \textquotedblleft{}the side of a river\textquotedblright{}.
\end{cousinweb}

\begin{grammarlens}[title={What you've built}]
One more word for getting around.
\end{grammarlens}

\subsection*{Your turn}
\begin{itemize}
\raggedright
  \item \emph{Say it:} \emph{ḍiffa}
  \item \emph{Say it:} \emph{ḍiffa}, once more
  \item \emph{Say it:} \emph{jalasnā ʿalā ḍiffat an-nahr.} — We sat on the riverbank.
  \item \emph{Recall:} say the Arabic for a river source, then the Arabic for a river mouth, then say \emph{ḍiffa} again
\end{itemize}

\begin{tcolorbox}[breakable,colback=teal!4,colframe=teal!35!black,title={Before you move on}]
What does \ar{ضفة} mean? (\textquotedblleft{}A riverbank\textquotedblright{}.) Say it once more.
\end{tcolorbox}
\section[\texorpdfstring{\ar{عذب} (ʿadhb)}{عذب (ʿadhb)}]{\ar{عذب} (ʿadhb) — sweet (of water); fresh}

\label{lesson:AR-C470-adhb}



\begin{quote}
Before the new one: say the Arabic for a river mouth, then the Arabic for a riverbank.
\end{quote}

\subsection*{You'll want to know: \ar{عذب}}
\textbf{\ar{عذب}} — \emph{ʿadhb} — \textquotedblleft{}sweet (of water); fresh\textquotedblright{}.

\begin{quote}
\ar{ماء} \ar{الينبوع} \ar{عذب}.
\end{quote}

\begin{quote}
\emph{māʾ al-yanbūʿ ʿadhb.}
\end{quote}

\begin{quote}
The water of the spring is sweet.
\end{quote}

Of water, sweet means fresh, not salty.

\begin{cousinweb}
\textbf{\ar{عذب}} (\emph{ʿadhb}) is built on the native Arabic root \emph{ʿ-dh-b}, \textquotedblleft{}to be sweet (of water)\textquotedblright{} (Lane).
\end{cousinweb}

\begin{grammarlens}[title={What you've built}]
One more word for getting around.
\end{grammarlens}

\subsection*{Your turn}
\begin{itemize}
\raggedright
  \item \emph{Say it:} \emph{ʿadhb}
  \item \emph{Say it:} \emph{ʿadhb}, once more
  \item \emph{Say it:} \emph{māʾ al-yanbūʿ ʿadhb.} — The water of the spring is sweet.
  \item \emph{Recall:} say the Arabic for a river mouth, then the Arabic for a riverbank, then say \emph{ʿadhb} again
\end{itemize}

\begin{tcolorbox}[breakable,colback=teal!4,colframe=teal!35!black,title={Before you move on}]
What does \ar{عذب} mean? (\textquotedblleft{}Sweet (of water); fresh\textquotedblright{}.) Say it once more.
\end{tcolorbox}
\section[(dialogue)]{Words from chapters 463–470 — a review}

\label{lesson:AR-R470-review-463-470}



\begin{quote}
Say the Arabic for climate and global warming.
\end{quote}

\subsection*{Your turn}
Say these aloud:

\begin{itemize}
\raggedright
  \item climate, global warming, the atmosphere, an emission, carbon
  \item pollution, polluted, waste, to pollute, to sort
  \item energy, gas, crude oil, coal, to consume
  \item fine dust, dusty, desertification, a sand dune, to choke
  \item a nature reserve, endangered, to protect, to die out, to squander
  \item plant life, a shrub, to wilt, to bloom, to bear fruit
  \item soil, fertiliser, a crop, fertile, to uproot
  \item a spring, a river source, a river mouth, a riverbank, sweet
\end{itemize}

\begin{tcolorbox}[breakable,colback=teal!4,colframe=teal!35!black,title={Before you move on}]
Climate? (\textbf{\ar{مناخ}}, \emph{munākh}.) Global warming? (\textbf{\ar{احتباس} \ar{حراري}}, \emph{iḥtibās ḥarārī}.) The atmosphere? (\textbf{\ar{غلاف} \ar{جوي}}, \emph{ghilāf jawwī}.) An emission? (\textbf{\ar{انبعاث}}, \emph{inbiʿāth}.) Carbon? (\textbf{\ar{كربون}}, \emph{karbūn}.) Pollution? (\textbf{\ar{تلوث}}, \emph{talawwuth}.) Polluted? (\textbf{\ar{ملوث}}, \emph{mulawwath}.) Waste? (\textbf{\ar{نفايات}}, \emph{nufāyāt}.) To pollute? (\textbf{\ar{لوث}}, \emph{lawwatha}.) To sort? (\textbf{\ar{فرز}}, \emph{faraza}.) Energy? (\textbf{\ar{طاقة}}, \emph{ṭāqa}.) Gas? (\textbf{\ar{غاز}}, \emph{ghāz}.) Crude oil? (\textbf{\ar{نفط}}, \emph{nafṭ}.) Coal? (\textbf{\ar{فحم}}, \emph{faḥm}.) To consume? (\textbf{\ar{استهلك}}, \emph{istahlaka}.) Fine dust? (\textbf{\ar{غبار}}, \emph{ghubār}.) Dusty? (\textbf{\ar{مغبر}}, \emph{mughbarr}.) Desertification? (\textbf{\ar{تصحر}}, \emph{taṣaḥḥur}.) A sand dune? (\textbf{\ar{كثيب}}, \emph{kathīb}.) To choke? (\textbf{\ar{اختنق}}, \emph{ikhtanaqa}.) A nature reserve? (\textbf{\ar{محمية}}, \emph{maḥmiyya}.) Endangered? (\textbf{\ar{مهدد}}, \emph{muhaddad}.) To protect? (\textbf{\ar{حافظ}}, \emph{ḥāfaẓa}.) To die out? (\textbf{\ar{انقرض}}, \emph{inqaraḍa}.) To squander? (\textbf{\ar{أهدر}}, \emph{ahdara}.) Plant life? (\textbf{\ar{نبات}}, \emph{nabāt}.) A shrub? (\textbf{\ar{شجيرة}}, \emph{shujayra}.) To wilt? (\textbf{\ar{ذبل}}, \emph{dhabala}.) To bloom? (\textbf{\ar{أزهر}}, \emph{azhara}.) To bear fruit? (\textbf{\ar{أثمر}}, \emph{athmara}.) Soil? (\textbf{\ar{تربة}}, \emph{turba}.) Fertiliser? (\textbf{\ar{سماد}}, \emph{samād}.) A crop? (\textbf{\ar{محصول}}, \emph{maḥṣūl}.) Fertile? (\textbf{\ar{خصب}}, \emph{khaṣib}.) To uproot? (\textbf{\ar{اقتلع}}, \emph{iqtalaʿa}.) A spring? (\textbf{\ar{ينبوع}}, \emph{yanbūʿ}.) A river source? (\textbf{\ar{منبع}}, \emph{manbaʿ}.) A river mouth? (\textbf{\ar{مصب}}, \emph{maṣabb}.) A riverbank? (\textbf{\ar{ضفة}}, \emph{ḍiffa}.) Sweet? (\textbf{\ar{عذب}}, \emph{ʿadhb}.)
\end{tcolorbox}
